\section{Onde aprofundar}

\subsection*{Documentação}
\small
\begin{longtable}{@{}p{5cm} p{7.2cm}@{}}
\toprule
Fonte & Para quê \\
\midrule \endhead
\citetitle{doc:quarkus-writing-extensions} & Build steps, build items,
  recorders, índice e as propriedades de depuração deste capítulo \\
\citetitle{doc:quarkus-cdi-reference} & Descoberta de beans, remoção de
  beans não usados, membros \texttt{private} e limitações em relação ao
  CDI \\
\citetitle{doc:quarkus-class-loading} & Os class loaders do dev mode e
  do pacote de produção (seção~\ref{sec:class-loading}) \\
\citetitle{doc:graalvm-native-basics} & Inicialização de classes em
  build time e o que vai para o binário nativo \\
\bottomrule
\end{longtable}
\normalsize

\subsection*{Ferramentas}
\citetitle{doc:jandex} (o indexador) e \citetitle{doc:gizmo} (o gerador
de bytecode). Para ler o que eles produzem, as mesmas ferramentas da
seção~\ref{sec:ferramentas}: \texttt{javap}, JFR e JITWatch.

Os livros sobre Quarkus estão no capítulo~\ref{cap:onde-aprender}.

\subsection*{Pesquisa acadêmica}
Acesso (aberto ou pago) e links no apêndice~\ref{ap:pesquisa}.
\begin{itemize}
\pesquisa{paper:native-image-initialize-once}{O paper do Native Image:
  análise em mundo fechado e inicialização de classes no build, a base do
  modo nativo do Quarkus.}
\pesquisa{paper:one-vm-to-rule-them-all}{O manifesto do Graal e do Truffle:
  uma VM para muitas linguagens.}
\end{itemize}
